\documentclass[12pt,a4paper]{program-proposal}
\usepackage[none]{hyphenat}
\usepackage{multirow}

\title{Coloring Cubes\\{\normalsize kaNita kAnakam 2026}}
\author{~}
\date{\today}
\begin{document}
\maketitle
\thispagestyle{empty}
A cube has six faces. If you have six different colours and use one colour on each face, how many different cubes can you make? To answer this, we first have to decide what ``different" means, and to help with that we will construct an origami cube.

\subsection{Activity 1:} Work in pairs. Each person makes 3 Sonobe units. Put the six units together to build an origami cube. As it's paper, we can mark/write on the faces of the cube.

\subsection{Activity 2:} One student closes their eyes while the other either rotates the cube, or leaves it as it was. The first student opens their eyes and has to tell whether the cube has been rotated or not. 

We then explore the different rotations of a cube and discover that there are 24 ways to rotate it without changing its overall structure. We informally explore what happens when two rotations are performed one after another.

At this point we will return to the original question. We can expect two ways of answering the question. One, the 6 faces can be painted in $6!$ ways. Each cube is counted 24 times, so the number of different ways = $6!/24$. Two, fix a face, pick the opposite face in 5 ways, the remaining 4 faces form a circular permutation in 6 ways, giving a total of $5\times 6$ ways. This bit in my experience takes the most amount of time. We can do problem 2 and more if we have time, if not, the rest will be homework.

\subsection{Problem 2:} Instead of 6 colors, what if we had only 2 colors and had to color 3 faces in one color and the other 3 in the other color. How many different cubes can be made like that?

\subsection{Problem 2:} What if we had only 3 colors and had to color 2 faces in each color. How many different cubes can be painted like that?

\subsection{Activity 3:} If time permits, create triangle units to make a tetrahedron. Repeat problem 1, 2, 3 and Activity 2.

\subsection{Homework:} Use the triangle units to create an icosahedron. Repeat problem 1, 2, 3 and Activity 2.
\end{document}

\begin{tabular}{ |p{0.23\textwidth}|p{0.66\textwidth}|  }
 \hline
 \textbf{Session, tentative date \& Theme} &\textbf{Outline}\\

 \multirow{3}{0.23\textwidth}{Session 4\\\\\emph{}}  & \\
\hline

% \hline
%  \multirow{2}{0.22\textwidth}{Session 8\\20 December 2026}  & Fractals and aperiodic tiling - II  & Iterated function system\\
% \hline
%  \multirow{2}{0.22\textwidth}{Session 9\\17 January 2027}  &  \\
% \hline
%  \multirow{2}{0.22\textwidth}{Session 10\\ }  & Chess and logic puzzles. Revisit the graph theory concepts that we developed early on. Application of graph theory\\
\hline
\end{tabular}




Following is a proposal for a sequence of mathematics workshops designed to introduce students to a range of mathematical ideas through accessible and exploratory entry points. The sessions aim to cultivate mathematical thinking by engaging students with visual, tactile, and pattern-based approaches to mathematics.

Although certain sessions build cumulatively on concepts developed earlier in the program, the structure also allows for partial participation, as not all sessions require prior attendance. Each session is designed to be of 150 minutes with a break of 15 minutes in between.

The program is organised around three broad thematic strands.

The first strand focuses on hands-on explorations through paper folding and origami-based activities conducted over one or two sessions. These sessions are intended to broaden students' perceptions of mathematics beyond computation and arithmetic, highlighting the role of spatial reasoning, visualisation, and embodied interaction in mathematical thinking. Through origami constructions, students will be introduced to geometric ideas such as polyhedra, alongside an initial exposure to elementary graph-theoretic concepts.

The second strand centres on symmetry as an intuitive entry point into abstract algebraic thinking, particularly group theory. Spanning three to four sessions, this component assumes familiarity with ideas introduced earlier in the program. Students will investigate wallpaper groups, frieze groups, and tessellations, while also drawing connections to the polyhedral models constructed during the origami sessions. The emphasis will be on recognising and analysing structural regularities through concrete and visual representations.

The subsequent sessions focus on mathematical patterns and periodicity. Building on observations emerging from the symmetry activities, students will examine different forms of repeating structures and develop an understanding of periodic behaviour in mathematical and real-world contexts. These sessions aim to leverage students' natural tendency to identify patterns in their surroundings as a foundation for mathematical abstraction. Subject to time availability, the program may conclude with explorations of modified chessboard problems that utilise introductory graph-theoretic ideas developed in the initial sessions.

The program is designed for students in the 13 to 17 age group (approximately Classes 8 to 10) and is intended for cohorts of around 30 participants.

Tentative dates: \date{19 July 2026}, \date{16 August 2026}, \date{30 August 2026}, \date{6 September 2026}, \date{20 September 2026}, \date{11 October 2026}, \date{18 October 2026}, \date{25 October 2026}, \date{ 1 November 2026}, \date{ 8 November 2026} 
\section{Origami}
\subsection{Outline}    
Mathematics is frequently perceived as a discipline mediated primarily through symbolic notation and formal manipulation on paper. This workshop seeks to broaden that perception by engaging participants in mathematical exploration through paper folding and origami-based inquiry. 

On the first set of the first set of instructions over will be constructing Bolera using modules using modular look at oil, carrots, characteristics, and such the second set would be on  flat-foldability: the study of conditions under which a crease pattern can be folded into a planar form without tearing or crumpling.

Using hands-on activities as a mode of mathematical investigation, participants will engage in processes of conjecture formation, empirical testing, and collective reasoning. Through guided exploration, the workshop will introduce participants to mathematical argumentation and proof in an accessible and experiential setting. No prior background in advanced mathematics or origami is assumed.

In addition to constructing origami models, participants will encounter mathematical ideas emerging from the folding process itself, illustrating how tactile and visual experiences can support the development of geometric and structural reasoning.

\section{Symmetry}
